\documentclass[12pt]{article}
\usepackage{amsmath,amssymb,amsfonts}
\usepackage[margin=1in]{geometry}
\usepackage{bm}

\newcommand{\vect}[1]{\mathbf{#1}}
\newcommand{\grad}{\boldsymbol{\nabla}}
\newcommand{\curl}{\grad\times}
\newcommand{\divr}{\grad\cdot}

\begin{document}

\section*{Problem 10}

\textbf{Problem:}
\begin{itemize}
\item[(a)] What class of physical phenomena are described by Helmholtz's equation, $\nabla^2\psi + k^2\psi = 0$?
\item[(b)] Interpret this equation physically in the light of the discussion of the physical meaning of the Laplacian, and square your interpretation with the answer to part (a).
\item[(c)] Derive the law of conservation of charge from Maxwell's equations.
\item[(d)] Express the electric field in terms of the scalar potential $\varphi$ and the vector potential $\vect{A}$ where $\vect{B} = \curl\vect{A}$. (\textit{Note:} For the general time-dependent case, $\vect{E}$ is not equal to $-\grad\varphi$ alone.)
\end{itemize}

\bigskip
\textbf{Solution:}

\subsection*{(a) Physical phenomena described by Helmholtz's equation}

The Helmholtz equation $\nabla^2\psi + k^2\psi = 0$ describes \textbf{time-harmonic (monochromatic) wave phenomena}.

It arises when we substitute a harmonic time dependence $\Psi(\vect{r},t) = \psi(\vect{r})\,e^{-i\omega t}$ into the wave equation:
\[
\nabla^2\Psi = \frac{1}{c^2}\frac{\partial^2\Psi}{\partial t^2} \quad \Longrightarrow \quad \nabla^2\psi\,e^{-i\omega t} = \frac{-\omega^2}{c^2}\,\psi\,e^{-i\omega t} \quad \Longrightarrow \quad \nabla^2\psi + k^2\psi = 0,
\]
where $k = \omega/c$ is the wave number.

Examples include:
\begin{itemize}
\item Electromagnetic waves (light, radio, etc.) at a single frequency
\item Sound waves and acoustic vibrations
\item Quantum mechanical wavefunctions of free particles (where $k^2 = 2mE/\hbar^2$)
\item Vibrations of membranes and cavities (eigenvalue problems)
\end{itemize}

\subsection*{(b) Physical interpretation via the Laplacian}

As discussed in Section 1.7 of Byron \& Fuller, the Laplacian $\nabla^2\psi$ at a point measures the difference between the average value of $\psi$ in a neighborhood and its value at that point (Eq.~1.71):
\[
\nabla^2\psi \propto \langle\psi\rangle_{\text{neighborhood}} - \psi(\vect{r}).
\]

The Helmholtz equation $\nabla^2\psi = -k^2\psi$ says:
\[
\langle\psi\rangle_{\text{neighborhood}} - \psi(\vect{r}) \propto -k^2\psi(\vect{r}).
\]

\textbf{Interpretation:} Where $\psi > 0$, the average value of $\psi$ in the neighborhood is \emph{less} than $\psi$ at the point (the Laplacian is negative), meaning $\psi$ has a local maximum tendency. Where $\psi < 0$, the reverse holds. This oscillatory character---where the function is always being ``pulled back'' from its extremes---is precisely the behavior of a \textbf{standing wave}. The parameter $k$ controls the spatial frequency of oscillation.

This squares with part (a): wave phenomena are characterized by spatial oscillation, and the Helmholtz equation encodes exactly this oscillatory behavior through the Laplacian.

\subsection*{(c) Conservation of charge from Maxwell's equations}

Maxwell's equations in differential form (CGS/Gaussian units) are:
\begin{align}
\divr\vect{D} &= 4\pi\rho, \label{eq:M1} \\
\divr\vect{B} &= 0, \label{eq:M2} \\
\curl\vect{E} &= -\frac{1}{c}\frac{\partial\vect{B}}{\partial t}, \label{eq:M3} \\
\curl\vect{H} &= \frac{4\pi}{c}\vect{J} + \frac{1}{c}\frac{\partial\vect{D}}{\partial t}. \label{eq:M4}
\end{align}

Take the divergence of Eq.~\eqref{eq:M4}:
\[
\divr(\curl\vect{H}) = \frac{4\pi}{c}\divr\vect{J} + \frac{1}{c}\frac{\partial}{\partial t}(\divr\vect{D}).
\]

The left side vanishes identically ($\divr(\curl\vect{H}) = 0$). Using Eq.~\eqref{eq:M1} for $\divr\vect{D}$:
\[
0 = \frac{4\pi}{c}\divr\vect{J} + \frac{1}{c}\frac{\partial}{\partial t}(4\pi\rho) = \frac{4\pi}{c}\left(\divr\vect{J} + \frac{\partial\rho}{\partial t}\right).
\]

Therefore:
\[
\boxed{\frac{\partial\rho}{\partial t} + \divr\vect{J} = 0}.
\]

This is the \textbf{continuity equation} expressing conservation of electric charge. It states that the rate of decrease of charge density at a point equals the divergence of current density (net outflow of current) at that point.

\textit{Note:} It was precisely to ensure this conservation law that Maxwell added the displacement current term $\frac{1}{c}\frac{\partial\vect{D}}{\partial t}$ to Amp\`ere's law.

\subsection*{(d) Electric field in terms of potentials}

Since $\divr\vect{B} = 0$, we can write $\vect{B} = \curl\vect{A}$ for some vector potential $\vect{A}$.

Substituting into Faraday's law, Eq.~\eqref{eq:M3}:
\[
\curl\vect{E} = -\frac{1}{c}\frac{\partial}{\partial t}(\curl\vect{A}) = -\curl\!\left(\frac{1}{c}\frac{\partial\vect{A}}{\partial t}\right).
\]

Rearranging:
\[
\curl\!\left(\vect{E} + \frac{1}{c}\frac{\partial\vect{A}}{\partial t}\right) = \vect{0}.
\]

Since the curl of this expression vanishes, it must be the gradient of a scalar (Eq.~1.83):
\[
\vect{E} + \frac{1}{c}\frac{\partial\vect{A}}{\partial t} = -\grad\varphi,
\]

where $\varphi$ is the scalar potential. Therefore:
\[
\boxed{\vect{E} = -\grad\varphi - \frac{1}{c}\frac{\partial\vect{A}}{\partial t}}.
\]

In the static case ($\partial\vect{A}/\partial t = 0$), this reduces to $\vect{E} = -\grad\varphi$. In the general time-dependent case, the additional term $-\frac{1}{c}\frac{\partial\vect{A}}{\partial t}$ is essential. (In SI units, the factor $1/c$ is absent: $\vect{E} = -\grad\varphi - \frac{\partial\vect{A}}{\partial t}$.) \hfill $\blacksquare$

\end{document}
